
\subsection{Funamentals of cardinality}
Regardless of the definition of cardinality that we choose to use, we denote $|A|$ the 
cardinality of $A$. We will propose two different definition for cardinality. We will use the 
notation $|A| \darr |B|$ to denote that $|A|$ and $|B|$ are in mutual injection, 
thus that both an injection from $A$ to $B$ and an injection from $B$ to $A$ exist. Bijection 
clearly implies mutual injection. The Bernstein's theorem (theorem \ref{thm:Bernstein}) states that 
if there is a mutual injection implies mutual bijection, but right 
now, we can't prove it yet. \\

\begin{definition}[Cardinality with AC]
\label{def:cardinalityAC}
The cardinality of a set $A$ is the smallest ordinal in bijection with $A$, 
Assuming well ordering is essential to this definition.
\end{definition}

\begin{definition}[Scott's cardinality]
\label{def:cardinalityAC}
The cardinality of a set $A$ is the set of all sets in mutal injection with $A$ that also have 
minimal rank. \\
\end{definition}

\begin{definition}[Cardinality comparison]
We can compare cardinalities using $\leq$:

$$ |A| \leq |B| \iff \exists f : A \rightarrow B \;\;\texttt{injective} $$
\end{definition}

\begin{definition}[Cardinal sum]
We can define the sum for cardinal arithmetic as the following operation, which 
is defined as the cardinality of the "disjoint union":

$$ |A| + |B| = | (\{0\} \times A) \cup (\{1\} \times B) | $$
\end{definition}

\begin{definition}[Cardinal product]
We can define the sum for cardinal arithmetic as the following operation, which 
is defined as the cardinality of the set product:

$$ |A| \times |B| = | A \times B | $$
\end{definition}


\begin{definition}[Cardinal exponentiation]
We can define the sum for cardinal arithmetic as the following operation, which 
is defined as the cardinality of the set of all functions having domain and codomain 
as the operands:

$$ |A|^{|B|} = | { f : B \rightarrow A } | $$
\end{definition}
    
    
Please be very aware that using the definition of cardinality of a given set as the 
smallest ordinal in bijection with such set, you are still supposed to use the 
cardinal arithmetic operations defined above. Ordinal arithmetic is not the same as
cardinal arithmetic, and, since not every ordinal is a cardinal, we can't use the
ordinals arithmetic operations on cardinals. For example $\omega \neq \omega + 1$ but
$|A| + 1 = |A|$ for every set $A$ in bijection with omega. \\

\newpage